% GENERATED FILE. Edit canonical lessons, then run npm run generate:books.
% canonical-source-hash: fnv1a64:d5cc5ef0eda6f838
% canonical-lessons: JA-C289-hanayome, JA-C289-hanamuko, JA-C289-chikau, JA-C289-shoutai-suru, JA-C289-fuufu

\chapter{A Wedding}
\label{ch:ja-a-wedding}
\hlchaptermodality{289}

\begin{chapteropening}
\textbf{By the end of this chapter:} \emph{I can say who married and what they promised.}

Complete the last lesson of chapter 289: I can say who married and what they promised.
\end{chapteropening}

\section[\texorpdfstring{\ja{はなよめ} (hanayome)}{はなよめ (hanayome)}]{\ja{はなよめ} (hanayome) — a bride}

\label{lesson:JA-C289-hanayome}



\begin{quote}
Before the new one: say the Japanese for lively, then the Japanese for a celebration.
\end{quote}

\subsection*{You'll want to know: \ja{はなよめ}}
\textbf{\ja{はなよめ}} — \emph{hanayome} — \textquotedblleft{}a bride\textquotedblright{}.

\begin{quote}
\textbf{\ja{はなよめは} \ja{しろい} \ja{きものを} \ja{きていました。}} — \emph{hanayome wa shiroi kimono o kite imashita} — \textquotedblleft{}The bride wore a white kimono.\textquotedblright{}
\end{quote}

\begin{cousinweb}
\textbf{\ja{はなよめ}} (\emph{hanayome}) is a native Japanese word (wago), inherited rather than borrowed from Chinese, built from \emph{hana}, \textquotedblleft{}flower\textquotedblright{}, and \emph{yome}, \textquotedblleft{}bride\textquotedblright{} (JMdict, KANJIDIC2).
\end{cousinweb}

\begin{grammarlens}[title={What you've built}]
One more word for telling what happened.
\end{grammarlens}

\subsection*{Your turn}
\begin{itemize}
\raggedright
  \item \emph{Say it:} \emph{hanayome}
  \item \emph{Say it:} \emph{hanayome}, once more
  \item \emph{Say it:} \emph{hanayome wa shiroi kimono o kite imashita} — The bride wore a white kimono.
  \item \emph{Recall:} say the Japanese for lively, then the Japanese for a celebration, then say \emph{hanayome} again
\end{itemize}

\begin{tcolorbox}[breakable,colback=teal!4,colframe=teal!35!black,title={Before you move on}]
What does \ja{はなよめ} mean? (\textquotedblleft{}A bride\textquotedblright{}.) Say it once more.
\end{tcolorbox}
\section[\texorpdfstring{\ja{はなむこ} (hanamuko)}{はなむこ (hanamuko)}]{\ja{はなむこ} (hanamuko) — a bridegroom}

\label{lesson:JA-C289-hanamuko}



\begin{quote}
Before the new one: say the Japanese for a celebration, then the Japanese for a bride.
\end{quote}

\subsection*{You'll want to know: \ja{はなむこ}}
\textbf{\ja{はなむこ}} — \emph{hanamuko} — \textquotedblleft{}a bridegroom\textquotedblright{}.

\begin{quote}
\textbf{\ja{はなむこは} \ja{わらっていました。}} — \emph{hanamuko wa waratte imashita} — \textquotedblleft{}The groom was smiling.\textquotedblright{}
\end{quote}

\begin{cousinweb}
\textbf{\ja{はなむこ}} (\emph{hanamuko}) is a native Japanese word (wago), inherited rather than borrowed from Chinese, built from \emph{hana}, \textquotedblleft{}flower\textquotedblright{}, and \emph{muko}, \textquotedblleft{}son-in-law\textquotedblright{} (JMdict, KANJIDIC2).
\end{cousinweb}

\begin{grammarlens}[title={What you've built}]
One more word for telling what happened.
\end{grammarlens}

\subsection*{Your turn}
\begin{itemize}
\raggedright
  \item \emph{Say it:} \emph{hanamuko}
  \item \emph{Say it:} \emph{hanamuko}, once more
  \item \emph{Say it:} \emph{hanamuko wa waratte imashita} — The groom was smiling.
  \item \emph{Recall:} say the Japanese for a celebration, then the Japanese for a bride, then say \emph{hanamuko} again
\end{itemize}

\begin{tcolorbox}[breakable,colback=teal!4,colframe=teal!35!black,title={Before you move on}]
What does \ja{はなむこ} mean? (\textquotedblleft{}A bridegroom\textquotedblright{}.) Say it once more.
\end{tcolorbox}
\section[\texorpdfstring{\ja{ちかう} (chikau)}{ちかう (chikau)}]{\ja{ちかう} (chikau) — to vow; to swear}

\label{lesson:JA-C289-chikau}



\begin{quote}
Before the new one: say the Japanese for a bride, then the Japanese for a bridegroom.
\end{quote}

\subsection*{You'll want to know: \ja{ちかう}}
\textbf{\ja{ちかう}} — \emph{chikau} — \textquotedblleft{}to vow; to swear\textquotedblright{}.

\begin{quote}
\textbf{\ja{ふたりは} \ja{じんじゃで} \ja{ちかいました。}} — \emph{futari wa jinja de chikaimashita} — \textquotedblleft{}The two made their vows at the shrine.\textquotedblright{}
\end{quote}

Polite \textbf{\ja{ちかいます}}, \emph{chikaimasu}; a godan verb. Past (polite): \textbf{\ja{ふたりが} \ja{ちかいました。}}, \emph{futari ga chikaimashita}.

\begin{cousinweb}
\textbf{\ja{ちかう}} (\emph{chikau}) is a native Japanese word (wago), inherited rather than borrowed from Chinese; its kanji was attached later as a kun'yomi (JMdict).
\end{cousinweb}

\begin{grammarlens}[title={What you've built}]
One more word for telling what happened.
\end{grammarlens}

\subsection*{Your turn}
\begin{itemize}
\raggedright
  \item \emph{Say it:} \emph{chikau}
  \item \emph{Say it:} \emph{chikau}, once more
  \item \emph{Say it:} \emph{futari wa jinja de chikaimashita} — The two made their vows at the shrine.
  \item \emph{Recall:} say the Japanese for a bride, then the Japanese for a bridegroom, then say \emph{chikau} again
\end{itemize}

\begin{tcolorbox}[breakable,colback=teal!4,colframe=teal!35!black,title={Before you move on}]
What does \ja{ちかう} mean? (\textquotedblleft{}To vow; to swear\textquotedblright{}.) Say it once more.
\end{tcolorbox}
\section[\texorpdfstring{\ja{しょうたいする} (shōtai suru)}{しょうたいする (shōtai suru)}]{\ja{しょうたいする} (shōtai suru) — to invite formally; to host}

\label{lesson:JA-C289-shoutai-suru}



\begin{quote}
Before the new one: say the Japanese for a bridegroom, then the Japanese for to vow.
\end{quote}

\subsection*{You'll want to know: \ja{しょうたいする}}
\textbf{\ja{しょうたいする}} — \emph{shōtai suru} — \textquotedblleft{}to invite formally; to host\textquotedblright{}.

\begin{quote}
\textbf{\ja{あには} \ja{しごとの} \ja{ともだちも} \ja{しょうたいしました。}} — \emph{ani wa shigoto no tomodachi mo shōtai shimashita} — \textquotedblleft{}My brother invited his friends from work as well.\textquotedblright{}
\end{quote}

Polite \textbf{\ja{しょうたいします}}, \emph{shōtaishimasu}; a \ja{する} verb. Past (polite): \textbf{\ja{あにが} \ja{ともだちを} \ja{しょうたいしました。}}, \emph{ani ga tomodachi o shōtai shimashita}.

\begin{cousinweb}
\textbf{\ja{しょうたいする}} (\emph{shōtai suru}) is Sino-Japanese (kango): two Chinese characters read with their on'yomi, readings borrowed from Chinese. They are \emph{shō}, \textquotedblleft{}invite\textquotedblright{} (Middle Chinese \emph{tsyew}) and \emph{tai}, \textquotedblleft{}wait\textquotedblright{} (\emph{dojX}) (KANJIDIC2; Middle Chinese in Baxter's transcription). The native verb \emph{suru}, \textquotedblleft{}to do\textquotedblright{}, makes it a verb.
\end{cousinweb}

\begin{grammarlens}[title={What you've built}]
One more word for telling what happened.
\end{grammarlens}

\subsection*{Your turn}
\begin{itemize}
\raggedright
  \item \emph{Say it:} \emph{shōtai suru}
  \item \emph{Say it:} \emph{shōtai suru}, once more
  \item \emph{Say it:} \emph{ani wa shigoto no tomodachi mo shōtai shimashita} — My brother invited his friends from work as well.
  \item \emph{Recall:} say the Japanese for a bridegroom, then the Japanese for to vow, then say \emph{shōtai suru} again
\end{itemize}

\begin{tcolorbox}[breakable,colback=teal!4,colframe=teal!35!black,title={Before you move on}]
What does \ja{しょうたいする} mean? (\textquotedblleft{}To invite formally; to host\textquotedblright{}.) Say it once more.
\end{tcolorbox}
\section[\texorpdfstring{\ja{ふうふ} (fūfu)}{ふうふ (fūfu)}]{\ja{ふうふ} (fūfu) — a married couple}

\label{lesson:JA-C289-fuufu}



\begin{quote}
Before the new one: say the Japanese for to vow, then the Japanese for to invite formally.
\end{quote}

\subsection*{You'll want to know: \ja{ふうふ}}
\textbf{\ja{ふうふ}} — \emph{fūfu} — \textquotedblleft{}a married couple\textquotedblright{}.

\begin{quote}
\textbf{\ja{となりの} \ja{ふうふは} \ja{わかいです。}} — \emph{tonari no fūfu wa wakai desu} — \textquotedblleft{}The couple next door are young.\textquotedblright{}
\end{quote}

\begin{cousinweb}
\textbf{\ja{ふうふ}} (\emph{fūfu}) is Sino-Japanese (kango): two Chinese characters read with their on'yomi, readings borrowed from Chinese. They are \emph{fū}, \textquotedblleft{}husband\textquotedblright{} (Middle Chinese \emph{pju}) and \emph{fu}, \textquotedblleft{}wife\textquotedblright{} (\emph{bjuwX}) (KANJIDIC2; Middle Chinese in Baxter's transcription).
\end{cousinweb}

\begin{grammarlens}[title={What you've built}]
One more word for telling what happened.
\end{grammarlens}

\subsection*{Your turn}
\begin{itemize}
\raggedright
  \item \emph{Say it:} \emph{fūfu}
  \item \emph{Say it:} \emph{fūfu}, once more
  \item \emph{Say it:} \emph{tonari no fūfu wa wakai desu} — The couple next door are young.
  \item \emph{Recall:} say the Japanese for to vow, then the Japanese for to invite formally, then say \emph{fūfu} again
\end{itemize}

\begin{tcolorbox}[breakable,colback=teal!4,colframe=teal!35!black,title={Before you move on}]
What does \ja{ふうふ} mean? (\textquotedblleft{}A married couple\textquotedblright{}.) Say it once more.
\end{tcolorbox}
